\documentclass[../../../include/open-logic-section]{subfiles}

\begin{document}

\olfileid{his}{set}{cantorplane}
\olsection{لکیر تے مستوی بارے کانتور}

شروڈر-برنسٹین دے ثبوت نال جڑے کجھ حالات دا تعلق اوس تاریخ نال اے جیہڑی اسیں \olref[his][set][pathology]{sec} وچ دَسی سی۔ چیتے کرو کہ 1877 وچ کانتور نے ثابت کیتا کہ مربع وچ عین اوہنے ای نقطے نیں جِنّے اوہدی اک بھجّا اُتے نیں۔ ایتھے اسیں اوہدے ثبوت دی پہلی کوشش پیش کراں گے۔

منّ لو $\unitline$ اکائی لکیر اے، یعنی نقطیاں دا سیٹ $[0,1]$۔ منّ لو $\unitsquare$ اکائی مربع اے، یعنی نقطیاں دا سیٹ $\unitline
\times \unitline$۔ ایہناں اصطلاحاں وچ کانتور نے ثابت کیتا کہ $\cardeq{\unitline}{\unitsquare}$۔ اوہنے ڈیڈیکائنڈ نوں اک سنیہا لکھیا جیہدے وچ بنیادی طور اُتے ایہ دلیل سی۔

\begin{thm}\ollabel{thm:cantorplane}
$\cardeq{\unitline}{\unitsquare}$
\end{thm}

\begin{proof}[ثبوت دا پہلا حصہ]
$a, b \in \unitline$ چُن لو۔ اوہناں نوں دوئی لکھت وچ لکھو۔ ایس طرح $0$ تے $1$ دیاں دو لامتناہی لَڑیاں، $a_1$، $a_2$، \dots، تے $b_1$، $b_2$، \dots، ملدیاں نیں، جیہڑیاں ایویں نیں:
\begin{align*}
a &= 0.a_1a_2a_3a_4\dots\\
b &= 0.b_1b_2b_3b_4\dots
\intertext{ہُن ایویں دِتّا فنکشن $f \colon \unitsquare \to \unitline$ لو:}
f(a, b) & = 0.a_1b_1a_2b_2a_3b_3a_4b_4\dots
\end{align*}
ہُن $f$ اک !!a{injection} اے، کیوں جے جے $f(a, b) = f(c,d)$، تاں $a_n =
c_n$ تے $b_n = d_n$ سارے $n \in \Nat$ لئی، سو $a = c$ تے $b =
d$۔
\end{proof}

بدقسمتی نال، جیویں ڈیڈیکائنڈ نے کانتور نوں دسیا، ایس نال اصل سوال دا جواب نہیں ملدا۔ $0.\dot{1}\dot{0} =
0.1010101010\ldots$ لو۔ سانوں چاہیدا اے کہ $f(a,b) = 0.\dot{1}\dot{0}$ ہووے، جتھے:
\begin{align*}
a&= 0.\dot{1}\dot{1} = 0.111111\ldots\\
b&= 0
\end{align*}
پر $a = 0.\dot{1}\dot{1} = 1$ اے۔ سو جدوں اسیں آکھدے آں کہ «$a$ تے $b$ نوں دوئی لکھت وچ لکھو»، تاں سانوں چُننا پیندا اے کہ \emph{کیہڑی} لکھت ورتنی اے۔ تے چونکہ $f$ اک \emph{فنکشن} اے، اسیں دوہاں ممکن لکھتاں وچوں بس \emph{اک} ای ورت سکدے آں۔ پر جے، مثال لئی، اسیں سوکھی لکھت ورت کے $a$ نوں «$1.000\ldots$» لکھیے، تاں کوئی ایہو جیہی جوڑی $\tuple{a, b}$ نہیں لبھدی جیہدے لئی $f(a, b) = 0.\dot{1}\dot{0}$ ہووے۔

مُکدی گل ایہ اے کہ ڈیڈیکائنڈ نے دسیا کہ کجھ دُہراندے اعشاری پھیلاؤ ممکن ہون کرکے کانتور دا فنکشن $f$ !!a{injection} اے پر !!a{surjection} \emph{نہیں}۔ سو کانتور نے بس $\cardle{\unitsquare}{\unitline}$ وکھایا اے تے $\cardeq{\unitsquare}{\unitline}$ \emph{نہیں} وکھایا۔

کانتور نے لگ بھگ جھٹ ڈیڈیکائنڈ نوں جواب لکھیا تے بنیادی طور اُتے ایہ دَسیا کہ ثبوت ایویں پورا ہو سکدا اے۔

\begin{proof}[پورا ثبوت]
سو اسیں وکھایا اے کہ $\cardle{\unitsquare}{\unitline}$۔ پر صاف اے کہ اک !!a{injection} اے جیہڑا $\unitline$ توں $\unitsquare$ ول جاندا اے۔ بس لکیر نوں مربع دی اک بھجّا دے نال سدھا رکھ دیو۔ سو $\cardle{\unitline}{\unitsquare}$ تے $\cardle{\unitsquare}{\unitline}$۔ شروڈر--برنسٹین دے قضیے نال
(\olref[sfr][siz][sb]{thm:schroder-bernstein})،
$\cardeq{\unitline}{\unitsquare}$۔
\end{proof}

پر ظاہر اے کہ کانتور آخری سطر ایس طرح پوری نہیں سی کر سکدا، کیوں جے شروڈر-برنسٹین دا قضیہ ہجے ثابت نہیں سی ہویا۔ بھاویں کانتور نے بعد وچ ایہنوں عام قیاس دی شکل دِتّی، ایہدا تسلّی بخش ثبوت 1897 تک نہیں ملیا۔ (ایس لئی کانتور نے 1877 وچ ای بعد وچ \olref{thm:cantorplane} دا اک ہور ثبوت دِتّا، جیہڑا شروڈر--برنسٹین اُتے نہیں سی ٹکدا۔)

\end{document}
